\documentclass[11pt,letterpaper]{article}
\usepackage[margin=1.05in,headheight=15pt]{geometry}
\usepackage{fontspec}
\setmainfont{Liberation Serif}
\setsansfont{Liberation Sans}
\setmonofont[Scale=0.81]{DejaVu Sans Mono}
\usepackage{amsmath,amssymb,amsthm,mathtools,mathrsfs}
\usepackage{microtype}
\usepackage{booktabs,tabularx,array,longtable}
\usepackage{enumitem}
\usepackage{xcolor}
\usepackage{hyperref}
\usepackage{fancyhdr}
\usepackage{listings}
\usepackage{caption}
\usepackage{needspace}
\usepackage{tikz}
\usetikzlibrary{arrows.meta,positioning}

\definecolor{ink}{HTML}{20354A}
\definecolor{accent}{HTML}{225B79}
\definecolor{muted}{HTML}{66717B}
\definecolor{lightpanel}{HTML}{F3F6F8}
\hypersetup{colorlinks=true,linkcolor=accent,urlcolor=accent,citecolor=accent,pdftitle={Golay Invariants and Arithmetic Realizations of M23},pdfauthor={Consolidated research note}}
\pagestyle{fancy}
\fancyhf{}
\fancyhead[L]{\small\sffamily\color{muted}MATHIEU GROUP \(M_{23}\)}
\fancyhead[R]{\small\sffamily\color{muted}GOLAY INVARIANTS \& INVERSE GALOIS THEORY}
\fancyfoot[C]{\small\sffamily\thepage}
\setlength{\parskip}{0.4em}
\setlength{\parindent}{0pt}
\setcounter{tocdepth}{2}
\setlist{itemsep=0.18em,topsep=0.25em}
\captionsetup{font=small,labelfont=bf}
\lstdefinestyle{compactpy}{language=Python,basicstyle=\ttfamily\scriptsize,keywordstyle=\color{accent}\bfseries,commentstyle=\color{muted},stringstyle=\color{ink},breaklines=true,columns=fullflexible,keepspaces=true,showstringspaces=false,frame=single,rulecolor=\color{muted!40},xleftmargin=0.5em,xrightmargin=0.5em,aboveskip=0.65em,belowskip=0.65em}
\newtheorem{theorem}{Theorem}[section]
\newtheorem{proposition}[theorem]{Proposition}
\newtheorem{lemma}[theorem]{Lemma}
\newtheorem{corollary}[theorem]{Corollary}
\theoremstyle{definition}
\newtheorem{remark}[theorem]{Remark}
\newcommand{\QQ}{\mathbb{Q}}
\newcommand{\FF}{\mathbb{F}}
\newcommand{\PP}{\mathbb{P}}
\newcommand{\Gal}{\operatorname{Gal}}
\newcommand{\Aut}{\operatorname{Aut}}
\newcommand{\Stab}{\operatorname{Stab}}
\newcommand{\Ni}{\operatorname{Ni}}
\newcommand{\Spec}{\operatorname{Spec}}
\newcommand{\Res}{\operatorname{Res}}
\newcommand{\tr}{\operatorname{Tr}}
\newcommand{\genus}{\operatorname{genus}}
\newcommand{\M}{M_{23}}
\newcommand{\G}{\mathcal{G}_{23}}
\newcommand{\HH}{\mathscr{H}}
\newcommand{\B}{\mathcal{B}_5}
\newcommand{\oct}{\mathcal{H}_7}

\begin{document}
\begin{titlepage}
\vspace*{0.85in}
{\sffamily\color{accent}\Large RESEARCH NOTE \hfill 14 NOVEMBER 2026}\par
\vspace{0.55in}
{\fontsize{27}{34}\selectfont\bfseries Golay-Code Invariants\\[0.04em] and the Construction of \(M_{23}\)-Polynomials}\par
\vspace{0.3in}
{\Large\color{muted}Connecting the 2026 inverse-Galois realization, real Nielsen-class structure, and explicit invariant theory}\par
\vspace{0.55in}
\noindent\color{accent}\rule{\linewidth}{1.4pt}\color{black}
\vspace{0.36in}
\begin{abstract}\noindent
The perfect binary Golay code \([23,12,7]\) has automorphism group \(M_{23}\), the last sporadic simple group realized as a Galois group over \(\QQ\) in the 2026 work of Huang, Jackson, Lee, Poonen, Pries and Zhang. A complementary real-structure analysis by Janik isolates three complex-conjugation-fixed covers in the seven-element Nielsen class \((2,23A,23B)\). This note synthesizes those results, incorporates the six-plus-one arithmetic orbit decomposition proved in the first paper, and builds two concrete invariants from the Steiner system \(S(4,7,23)\): a degree-five polynomial with 5,313 monomials and a degree-seven heptad polynomial with 253 monomials. The gradient of the latter supplies a 23-component \(M_{23}\)-equivariant covariant and hence a rigorous family of code-derived degree-23 polynomials over \(\QQ(t,u)\). We distinguish such new defining equations from genuinely new regular \(M_{23}\)-covers, establish an irreducible degree-253 heptad resolvent, and formulate a computational program for invariant-guided Hurwitz-space constructions. No new independent regular cover is claimed.
\end{abstract}
\vfill
\noindent\textbf{Source basis.} \cite{HJLPPZ26} (arXiv:2608.08538v2, 7 October 2026) and \cite{Janik26} (7 October 2026); computational data from the companion Golay invariant script included in Appendix~\ref{app:code}.\par
\vspace{0.2in}
\noindent\textbf{Status key.} We distinguish published theorems, elementary deductions proved here, computational checks, and proposed research tasks.\par
\end{titlepage}

\tableofcontents
\clearpage

\section{Question, scope, and summary}
The inverse Galois problem asks whether a finite group \(G\) occurs as \(\Gal(K/\QQ)\) for some finite Galois extension \(K/\QQ\). For \(M_{23}\), the existence question is now settled: a regular \(M_{23}\)-extension of \(\QQ(t)\) exists \cite[Theorem 1.3]{HJLPPZ26}. The question pursued here is more constructive:
\begin{quote}
Can the invariant theory of the perfect Golay code \emph{produce} algebraic equations with Galois group \(M_{23}\), rather than merely certify the group of previously constructed equations? Can it ultimately produce a distinct regular cover over \(\QQ\)?
\end{quote}

Three conclusions are separated throughout:
\begin{enumerate}[label=\textbf{\arabic*.}]
\item \textbf{Yes, for code-derived equations.} A Golay heptad polynomial yields an explicit equivariant root transformation, and with any known regular realization it gives a family of degree-23 equations with full \(M_{23}\) Galois group (Theorem~\ref{thm:transform}).
\item \textbf{No, not by root reparameterization alone.} Those equations describe the \emph{same} regular extension as the input. An alternative degree-253 heptad quotient changes the degree of the non-Galois cover but retains the same Galois closure.
\item \textbf{Still a research program for independent covers.} To obtain a different regular \(M_{23}\)-extension by invariant theory, one needs an independent rational Hurwitz point or a suitable arithmetic specialization of the universal invariant field, not just a new expression in an existing set of roots.
\end{enumerate}

\section{The binary Golay code, heptads, and \(M_{23}\)}
Let \(\G\subset\FF_2^{23}\) denote the perfect binary Golay code. It is the \([23,12,7]\) code obtained by puncturing the extended \([24,12,8]\) binary Golay code. Its automorphism group is
\[
\Aut(\G)\cong M_{23},\qquad |M_{23}|=10\,200\,960=2^7\cdot3^2\cdot5\cdot7\cdot11\cdot23.
\]
The extended code has automorphism group \(M_{24}\); the stabilizer of a coordinate of \(M_{24}\) is \(M_{23}\). The group \(M_{23}\) is 4-transitive in its natural 23-point action.

The weight distribution of \(\G\) is
\begin{center}
\begin{tabular}{@{}rrrrrrrrr@{}}\toprule
weight & 0 & 7 & 8 & 11 & 12 & 15 & 16 & 23\\\midrule
count & 1 & 253 & 506 & 1288 & 1288 & 506 & 253 & 1\\\bottomrule
\end{tabular}
\end{center}
A weight-seven word has a support of seven positions, called a \emph{heptad}. Let \(\oct\) be the family of 253 heptads. It forms the unique Steiner system \(S(4,7,23)\): every four-subset of \(\{1,\ldots,23\}\) lies in exactly one heptad. Double counting gives
\[
|\oct|=\frac{\binom{23}{4}}{\binom74}=253,\quad
\#\{H\in\oct:i\in H\}=\frac{\binom{22}{3}}{\binom63}=77.
\]
The code is perfect because
\[
2^{12}\sum_{j=0}^3\binom{23}{j}=2^{23},
\]
so radius-three Hamming balls centered on the codewords partition \(\FF_2^{23}\).

\section{What the two 2026 papers prove}
\subsection{A regular \(M_{23}\)-realization}
Huang, Jackson, Lee, Poonen, Pries and Zhang \cite{HJLPPZ26} construct a regular Galois extension \(E/\QQ(t)\) with
\[
\Gal(E/\QQ(t))\cong M_{23}.
\]
Their relevant geometric cover has branch locus
\[
S=\{0,\sqrt{-23},-\sqrt{-23}\},
\]
and inertia classes \((2,23A,23B)\). The degree-23 quotient by a point stabilizer \(M_{22}\) has genus four. Indeed, the involution has cycle type \(1^7 2^8\), and each order-23 element acts as a 23-cycle; Riemann--Hurwitz gives
\[
2g-2=-2(23)+8(2-1)+2(23-1)=6,\qquad g=4.
\]
The paper recovers a canonical genus-four curve numerically, descends a suitable cover to \(\QQ\), and produces a bidegree \((4,23)\) equation \(F(t,x)=0\). It then proves the geometric and arithmetic monodromy groups are \(M_{23}\). Its Appendix A uses the action on five-subsets to distinguish \(M_{23}\) from \(A_{23}\); this is a direct use of the Golay design, discussed in Section~\ref{sec:five}.

\subsection{Seven Nielsen classes and three real points}
Let \(\HH\) be the zero-dimensional Hurwitz scheme associated with the fixed three-point passport. Its geometric points correspond to the seven elements of the inner Nielsen class
\[
\Ni^{\mathrm{in}}(2,23A,23B),\qquad |\Ni^{\mathrm{in}}|=7.
\]
Janik \cite{Janik26} analyzes the action of complex conjugation on these seven covers by passing to braid orbits. With its stated conventions, the real involution is induced by
\[
(g_1,g_2,g_3)\longmapsto(g_3^{-1},g_2^{-1},g_1^{-1}),
\]
and acts on the seven orbit labels as \((2\;6)(3\;5)\), fixing \(\{0,1,4\}\). (The tuple form depends on the geometric-basis and permutation conventions; the induced number of fixed covers is invariant.) Thus
\[
\#\HH(\mathbb R)=3.
\]
The parity argument is elementary: complex conjugation is an involution of an odd set and hence has an odd number of fixed points. It guarantees a real point, not a rational point \cite[\S\S3--5]{Janik26}.

\subsection{The full arithmetic orbit structure: one plus six}
The second paper presents the remaining arithmetic orbit partition as a question. The first paper's revised Section 3.7 supplies the answer \cite[\S3.7, Proposition 3.14]{HJLPPZ26}:
\begin{equation}\label{eq:etale}
\HH\simeq \Spec\QQ\;\sqcup\;\Spec L,
\qquad
L=\QQ[z]/\bigl(z^6-z^5-3z^4-3z^3+z^2+5z+4\bigr).
\end{equation}
Here \(L\) is a degree-six field with signature \((2,2)\), whose Galois closure has group \(S_6\). The authors certify that a second geometric cover is defined over \(L\) but not over \(\QQ\), including a point-count obstruction at primes above 53. Since \(\deg\HH=7\), this leaves exactly one rational point and one six-element arithmetic orbit.

\begin{proposition}[Consequences for this fixed passport]
The fixed passport \((2,23A,23B)\) has exactly one rational cover and three real covers. The other two complex-conjugation-fixed covers belong to the single arithmetic orbit of size six and do not descend to \(\QQ\) as points of \(\HH\).
\end{proposition}
\begin{proof}
The decomposition \eqref{eq:etale} has exactly one \(\QQ\)-point, and its real points consist of the rational embedding and the two real embeddings of \(L\). This matches Janik's three fixed Nielsen orbits.
\end{proof}
This rules out finding a second inequivalent rational cover \emph{within the same fixed three-point passport}. It does not rule out another passport, a different branch locus, or another construction of a regular \(M_{23}\)-extension.

\section{The first exceptional invariant occurs in degree five}\label{sec:five}
Let \(x=(x_1,\ldots,x_{23})\) be independent variables over \(\QQ\), and define
\[
\B=\left\{S\subseteq[23]: |S|=5,\quad S\subseteq H\text{ for some }H\in\oct\right\}.
\]
No five-subset belongs to two distinct heptads, because any two heptads share at most three coordinates. Hence
\[
|\B|=253\binom75=5\,313.
\]
Define the squarefree invariant
\begin{equation}\label{eq:J5}
\boxed{\displaystyle J_5(x)=\sum_{S\in\B}\prod_{i\in S}x_i.}
\end{equation}
It contains 5,313 monomials, as opposed to \(\binom{23}{5}=33\,649\) possible squarefree degree-five monomials.

\begin{proposition}\label{prop:stabilizer}
The stabilizer of \(J_5\) in \(S_{23}\) is exactly \(M_{23}\).
\end{proposition}
\begin{proof}
Every automorphism of the Steiner system permutes the monomials of \(J_5\), so \(M_{23}\) fixes \(J_5\). Conversely, equality of squarefree polynomials implies any permutation fixing \(J_5\) preserves \(\B\). A seven-subset \(H\) is a heptad if and only if all its five-subsets lie in \(\B\): fixing any four-subset \(T\subset H\), each \(T\cup\{i\}\), for \(i\in H\setminus T\), belongs to a heptad, necessarily the unique heptad containing \(T\). Thus all elements of \(H\) lie in that one heptad. Consequently the permutation preserves \(\oct\), so it belongs to \(\Aut(S(4,7,23))=M_{23}\).
\end{proof}

Because \(M_{23}\) is 4-transitive, its orbits on monomials of total degree at most four agree with those of \(S_{23}\): an exponent pattern of degree at most four involves at most four distinct indices. Therefore all \(M_{23}\)-invariants of degree at most four are symmetric polynomials. Proposition~\ref{prop:stabilizer} shows that degree five is the first total degree with a non-symmetric \(M_{23}\)-invariant.

This mirrors a key step in \cite[Appendix A, Lemma A.5]{HJLPPZ26}: \(M_{23}\) has two orbits on five-subsets, while \(A_{23}\) has one. Their virtual character \(\eta=\chi_5-\chi_4\) detects this difference arithmetically over a finite field. The present proposal repurposes the same distinction to define polynomial invariants rather than merely test monodromy.

\section{The heptad polynomial and an equivariant covariant}
The more economical invariant is
\begin{equation}\label{eq:J7}
\boxed{\displaystyle J_7(x)=\sum_{H\in\oct}\prod_{i\in H}x_i.}
\end{equation}
It has 253 degree-seven monomials. Its stabilizer in \(S_{23}\) is also precisely \(M_{23}\), directly because the support of \(J_7\) is the set of heptads.

For \(1\le i\le23\), define
\begin{equation}\label{eq:D}
D_i(x)=\frac{\partial J_7}{\partial x_i}
=\sum_{\substack{H\in\oct\\i\in H}}\prod_{j\in H\setminus\{i\}}x_j.
\end{equation}
Each \(D_i\) is a sum of exactly 77 degree-six monomials. Set \(D(x)=(D_1(x),\ldots,D_{23}(x))\).

\begin{lemma}[Gradient equivariance]\label{lem:equiv}
For every \(g\in M_{23}\), with \(g\) acting by coordinate permutation,
\[
D(gx)=gD(x).
\]
\end{lemma}
\begin{proof}
Since \(J_7(gx)=J_7(x)\) and the coordinate action is by permutation matrices, the chain rule yields the asserted equivariance. Equivalently, the heptads containing \(i\) are carried bijectively to those containing \(g(i)\).
\end{proof}
The significance is that \(D\) produces a \emph{new 23-component root-like object}, rather than just a scalar invariant.

\section{A rigorous Golay-derived degree-23 polynomial family}
Let \(k=\QQ(t)\) and let \(f(t,X)\in k[X]\) be monic, separable and irreducible of degree 23, with regular splitting field \(E/k\) and Galois group identified with the natural \(M_{23}\)-permutation action. Label its roots \(\alpha_1,\ldots,\alpha_{23}\) compatibly with \(\oct\). This compatible labeling must be supplied explicitly in any implementation; an arbitrary root ordering will not work.

For an indeterminate \(u\), set
\begin{equation}\label{eq:beta}
\beta_i(t,u)=\alpha_i(t)+uD_i(\alpha_1(t),\ldots,\alpha_{23}(t))
\end{equation}
and define
\begin{equation}\label{eq:newpoly}
\boxed{\displaystyle F_{\mathrm{Golay}}(t,u,Y)=\prod_{i=1}^{23}(Y-\beta_i(t,u)).}
\end{equation}

\begin{theorem}[Code-derived polynomial family]\label{thm:transform}
The polynomial \(F_{\mathrm{Golay}}\) lies in \(\QQ(t)[u,Y]\) and has Galois group \(M_{23}\) over \(\QQ(t,u)\). For all but finitely many constants \(c\in\QQ\), the specialization \(F_{\mathrm{Golay}}(t,c,Y)\) is separable of degree 23 and has the same splitting field \(E\) over \(\QQ(t)\).
\end{theorem}
\begin{proof}
For \(\sigma\in\Gal(E/k)\), equivariance (Lemma~\ref{lem:equiv}) implies \(\sigma(\beta_i)=\beta_{\sigma(i)}\). The elementary symmetric functions in the \(\beta_i\) are therefore fixed by the Galois group and belong to \(k[u]\).

For distinct \(i,j\), the difference
\[
\beta_i-\beta_j=(\alpha_i-\alpha_j)+u(D_i(\alpha)-D_j(\alpha))
\]
is a nonzero polynomial in \(u\), because its constant coefficient \(\alpha_i-\alpha_j\) is nonzero. Hence the 23 transformed roots are pairwise distinct over \(E(u)\). The faithful permutation action on these roots identifies their splitting field over \(k(u)\) with \(E(u)\), yielding Galois group \(M_{23}\). All possible collisions under \(u=c\) lie among the finitely many zeros of the product of these nonzero linear polynomials; outside that finite set, the same faithful-action argument applies over \(k\).
\end{proof}

\begin{remark}[What is and is not new]
This theorem builds new \emph{defining equations}, using the code, from an existing regular extension. It does not independently prove a previously unknown regular extension exists, nor does it produce a distinct point of the fixed three-point Hurwitz scheme.
\end{remark}

\section{The Tschirnhaus obstruction and an exact resultant formula}
Let \(M_{22}=\Stab_{M_{23}}(1)\). Since \(D_1(\alpha)\) is fixed by \(M_{22}\),
\[
D_1(\alpha)\in E^{M_{22}}=k(\alpha_1).
\]
As \(k(\alpha_1)/k\) has degree 23, there is a unique representative \(P(t,X)\in k[X]\) of degree at most 22 with
\[
P(t,\alpha_1)=D_1(\alpha).
\]
Equivariance then ensures \(P(t,\alpha_i)=D_i(\alpha)\) for every \(i\). Therefore
\begin{equation}\label{eq:resultant}
\boxed{F_{\mathrm{Golay}}(t,u,Y)
=\Res_X\bigl(f(t,X),Y-X-uP(t,X)\bigr).}
\end{equation}
For generic \(u\), \(k(\beta_1)=k(\alpha_1)\), since the index of the stabilizer of \(\beta_1\) equals 23. Equation~\eqref{eq:resultant} is thus a code-defined \emph{Tschirnhaus transformation}, not an independent degree-23 covering. This is a principled obstruction to attempts that only reparameterize existing roots.

\section{An irreducible heptad resolvent and a second quotient cover}
The heptad incidence structure also gives a different natural permutation action of \(M_{23}\), on the 253 heptads. Let
\[
s_H=\sum_{i\in H}\alpha_i,\qquad
R_{253}(t,Z)=\prod_{H\in\oct}(Z-s_H).
\]
Galois invariance shows \(R_{253}\in k[Z]\). In this natural \(M_{23}\)-extension, the 253 sums are in fact automatically distinct, so no genericity hypothesis is needed.

\begin{lemma}[Distinct heptad sums]\label{lem:sums}
Assume characteristic zero and the full natural \(M_{23}\) Galois action on distinct \(\alpha_i\). If \(H\ne H'\) are heptads, then \(s_H\ne s_{H'}\).
\end{lemma}
\begin{proof}
Consider the \(M_{23}\)-equivariant linear map \(\phi:k^{23}\to E\) defined by \(\phi(c_1,\ldots,c_{23})=\sum_i c_i\alpha_i\). Since \(M_{23}\) is 2-transitive, its characteristic-zero permutation module decomposes into the trivial line and an irreducible 22-dimensional augmentation submodule \(V_0=\{c:\sum_i c_i=0\}\). The kernel of \(\phi\) is \(M_{23}\)-stable. It cannot contain all of \(V_0\), since that would force \(\alpha_i-\alpha_j=0\) for all \(i,j\). Hence \(\ker\phi\cap V_0=0\). The difference of the two incidence vectors \(1_H-1_{H'}\) is nonzero and lies in \(V_0\), so \(s_H-s_{H'}=\phi(1_H-1_{H'})\ne0\).
\end{proof}

\begin{proposition}\label{prop:res}
The polynomial \(R_{253}(t,Z)\) is irreducible of degree 253 over \(k\), and its splitting field is \(E\).
\end{proposition}
\begin{proof}
The group \(M_{23}\) acts transitively on the 253 heptads and, by Lemma~\ref{lem:sums}, acts on their sums with the same stabilizers. Thus \(R_{253}\) is an irreducible orbit polynomial of degree 253. The action on heptads is faithful because \(M_{23}\) is simple and acts nontrivially. Hence the splitting field has kernel trivial in \(\Gal(E/k)\), and must equal \(E\).
\end{proof}

Write \(Y\to\PP^1\) for the Galois closure. The point quotient \(Y/M_{22}\) is a genus-four degree-23 cover, whereas the heptad quotient \(Y/\Stab(H)\) has degree 253 and genus 46; the heptad stabilizer is the index-253 subgroup \(2^4\!:\!A_7\), appearing in the maximal-subgroup table of \cite[Table 2]{HJLPPZ26}. The underlying Galois closure is the same \(Y\).

The 23-point and 253-heptad sets admit an incidence correspondence with
\[
23\cdot77=253\cdot7=1771
\]
incident pairs. This correspondence can be represented by intermediate quotient curves and may be exploited for explicit trace, norm and resolvent computations. \emph{It is not itself a new regular \(M_{23}\)-extension.}

\section{The invariant-field route to independent arithmetic realizations}
There is a universal construction that requires no prior polynomial with Galois group \(M_{23}\). Take independent variables \(x_1,\ldots,x_{23}\) and let
\[
K=\QQ(x_1,\ldots,x_{23})^{M_{23}}.
\]
Then
\[
\prod_{i=1}^{23}(Y-x_i)\in K[Y],\qquad
\Gal\bigl(\QQ(x_1,\ldots,x_{23})/K\bigr)=M_{23}.
\]
This is a bona fide invariant-theoretic construction, but \(K\) has transcendence degree 23 over \(\QQ\), not one. The subtle task is to find arithmetic specializations that restrict this universal family to a regular extension over \(\QQ(t)\).

In geometric terms, one seeks appropriate rational curves or rational maps into the quotient \(\mathbb A^{23}/M_{23}\), with generic pullback retaining full monodromy \(M_{23}\). Neither the existence of \(J_5\) nor of \(J_7\) guarantees such a specialization. The real Nielsen-class analysis supplies an efficient combinatorial \emph{filter} for candidate Hurwitz points, but cannot by itself solve descent from \(\mathbb R\) to \(\QQ\).

\section{Computational program and proposed milestones}
\begin{enumerate}[label=\textbf{Phase \arabic*:},leftmargin=*]
\item \textbf{Establish the combinatorial model.} Generate the \([23,12,7]\) code, all 253 heptads, the 5,313 five-subsets in heptads, and the 23 sparse derivative expressions \(D_i\). Verify the Steiner property exactly. The included Appendix~\ref{app:code} is a reference implementation of this phase.
\item \textbf{Align the 23-point permutation actions.} Construct an explicit conjugating permutation between the \(M_{23}\) root action from \cite{HJLPPZ26} and the Golay-code coordinate action. Without this step, the code-based substitution can fail to be Galois-equivariant.
\item \textbf{Compute a concrete transformed equation.} Work in the degree-23 point-stabilizer field, recover \(P(t,X)\) from the heptad covariant and calculate the resultant \eqref{eq:resultant}. This avoids explicit arithmetic in the full splitting field of degree \(10\,200\,960\).
\item \textbf{Certify polynomial properties.} Verify coefficients lie in \(\QQ(t,u)\); test separability and irreducibility; establish the same splitting field as the input; optionally compute specializations \(u=c\in\QQ\) and optimize coefficient height.
\item \textbf{Compute the heptad quotient.} Construct \(R_{253}\), confirm its irreducibility and derive the incidence correspondence to the degree-23 quotient. Investigate whether it gives practical new resolvent or descent coordinates.
\item \textbf{Search beyond the fixed passport.} Enumerate other rationally admissible \(M_{23}\) branch-cycle passports, use braid-orbit and real-involution data to shortlist candidates, then impose low-degree Golay invariants as algebraic constraints in a Hurwitz-space computation. Use numerical or modular reconstruction, rational descent and independent monodromy certification before claiming a new regular cover.
\end{enumerate}

\subsection{What is already verified computationally}
The accompanying Python script was executed using the cyclic generator polynomial
\[
g(z)=z^{11}+z^9+z^7+z^6+z^5+z+1.
\]
It verified the weight distribution in Section 2, counted 253 heptads, checked that all \(\binom{23}{4}=8\,855\) four-subsets occur exactly once, counted \(5\,313\) distinguished five-subsets and 77 monomials in each component of \(D\). These computations validate the code data; they do not constitute a numeric construction of a new regular extension.

\section{Conclusions and research claims ledger}
\begin{center}
\small
\begin{tabularx}{\linewidth}{@{}p{0.45\linewidth}p{0.2\linewidth}X@{}}\toprule
Claim & Status & Basis\\\midrule
A regular \(M_{23}\)-extension over \(\QQ(t)\) exists & Published theorem & \cite{HJLPPZ26}\\
Seven geometric covers, three real fixed points & Published results & \cite{HJLPPZ26,Janik26}\\
Exactly one rational point in the fixed three-point passport & Published theorem & \cite[\S3.7]{HJLPPZ26}\\
\(\Stab_{S_{23}}(J_5)=M_{23}\) & Deduction proved here & Steiner-system reconstruction\\
Golay-gradient family has Galois group \(M_{23}\) & Deduction proved here & Theorem~\ref{thm:transform}\\
Degree-253 heptad resolvent is irreducible & Deduction proved here & Proposition~\ref{prop:res}\\
Constructing an independent regular cover from invariants & Open research target & Requires new arithmetic/descent\\\bottomrule
\end{tabularx}
\end{center}

The strongest immediate program is to use \(J_7\) to compute an explicit, exact Tschirnhaus transform, thereby demonstrating an effective code-to-equation mechanism. The higher-value, nontrivial objective is then to use \(J_5\) and related invariants as coordinates or constraints on \emph{other} Hurwitz spaces, where new rational points could yield genuinely independent regular \(M_{23}\)-realizations.

\appendix
\section{An explicit degree-23 arithmetic example}\label{app:example}
Example 1.2 of \cite{HJLPPZ26} gives the following polynomial with Galois group \(M_{23}\) over \(\QQ\), unramified outside \(\{2,3,23\}\):
{\small
\begin{align*}
p(X)={}&X^{23}-184X^{21}-1150X^{20}+26151X^{19}+18400X^{18}\\
&-1808490X^{17}+1545462X^{16}+67672923X^{15}\\
&-42732528X^{14}-1333395744X^{13}+290615166X^{12}\\
&+10550424369X^{11}+3700476348X^{10}+35123826654X^9\\
&-194398310718X^8-1023887308293X^7+3961650395556X^6\\
&+1949980486716X^5-28142323927002X^4+53599151839311X^3\\
&-46185312415788X^2+19169943578802X-3150159884154.
\end{align*}}
This example is a \emph{number-field specialization}; it should not be confused with the bivariate regular cover \(F(t,X)\). For the code-derived transformation in Section 6, the bivariate equation is the correct input if the goal is to retain a regular \(\QQ(t)\)-extension. Either input requires identifying the correct compatible root labeling before applying \(D\).

\section{Companion computational script}\label{app:code}
This self-contained script constructs the \([23,12,7]\) code and all sparse invariant supports. The root transformation routine requires correctly labeled roots in an arithmetic field; it does not perform the labeling or splitting-field computation for the user. The generated JSON file is auxiliary data, not necessary to compile this document.

\begin{lstlisting}[style=compactpy]
#!/usr/bin/env python3
"""Exact Golay/Steiner invariant data for a proposed M23 inverse-Galois experiment.

This script constructs the cyclic perfect binary Golay [23,12,7] code from
  g(x) = x^11 + x^9 + x^7 + x^6 + x^5 + x + 1 (0xAE3).
It then constructs its 253 weight-seven supports (heptads), the first
non-symmetric invariant J5 (supported on 5313 five-subsets), the heptad
invariant J7, and the 23 coordinates of the gradient of J7.

The abstract code data is *not* a new number-field realization. To use the
covariant root transform, one must supply roots already labelled so that
Galois action preserves these heptads (i.e. the action is the specific
coordinate copy of M23 chosen here).

Usage:
  python3 m23_golay_covariants.py [--json m23_golay_covariants.json]

The function derived_roots(alpha, u) performs the exact code-based root
transformation beta_i=alpha_i+u*d(J7)/d(x_i)(alpha).
The caller may use arithmetic elements in a suitably computed splitting
field to obtain polynomial coefficients via monic_coefficients(beta).
"""
from __future__ import annotations

import argparse
import itertools
import json
from collections import Counter
from math import comb
from pathlib import Path

N = 23
GOLAY_G = 0xAE3


def golay_codeword(message: int) -> int:
    assert 0 <= message < (1 << 12)
    result = 0
    for i in range(12):
        if message & (1 << i):
            result ^= GOLAY_G << i
    return result


def heptads() -> list[tuple[int, ...]]:
    blocks = [tuple(i for i in range(N) if word & (1 << i))
              for message in range(1 << 12)
              if (word := golay_codeword(message)).bit_count() == 7]
    return sorted(blocks)


def five_subsets(blocks: list[tuple[int, ...]]) -> list[tuple[int, ...]]:
    return sorted({s for h in blocks for s in itertools.combinations(h, 5)})


def gradient_monomials(blocks: list[tuple[int, ...]]) -> list[list[tuple[int, ...]]]:
    return [[tuple(j for j in h if j != i) for h in blocks if i in h]
            for i in range(N)]


def eval_monomial_sums(alpha, monomial_lists):
    """Evaluate a list of sparse monomial sums on elements of a commutative ring."""
    return [sum((prod_of(alpha[i] for i in monomial) for monomial in monomials), 0)
            for monomials in monomial_lists]


def prod_of(values):
    out = 1
    for x in values:
        out *= x
    return out


def derived_roots(alpha, u, blocks=None):
    """Return beta_i = alpha_i + u * dJ7/dx_i evaluated at alpha.

    Input alpha must be a list of exactly 23 correctly labelled root elements.
    The return values are algebraic elements, not yet rational coefficients.
    """
    if len(alpha) != N:
        raise ValueError('Expected exactly 23 labelled roots')
    if blocks is None:
        blocks = heptads()
    derivatives = eval_monomial_sums(alpha, gradient_monomials(blocks))
    return [alpha[i] + u * derivatives[i] for i in range(N)]


def monic_coefficients(roots):
    """Return ascending coefficients of prod(Y-root) in the input ring."""
    result = [1]
    for root in roots:
        next_coeffs = [0] * (len(result) + 1)
        for degree, coefficient in enumerate(result):
            next_coeffs[degree] -= root * coefficient
            next_coeffs[degree + 1] += coefficient
        result = next_coeffs
    return result


def verify(blocks):
    assert len(blocks) == 253 and len(set(blocks)) == 253
    wt = Counter(golay_codeword(i).bit_count() for i in range(1 << 12))
    assert wt == {0: 1, 7: 253, 8: 506, 11: 1288, 12: 1288,
                  15: 506, 16: 253, 23: 1}
    four_counts = Counter(s for h in blocks for s in itertools.combinations(h, 4))
    assert len(four_counts) == comb(23, 4) == 8855
    assert set(four_counts.values()) == {1}  # S(4,7,23)
    B5 = five_subsets(blocks)
    assert len(B5) == 253 * comb(7, 5) == 5313
    assert len(B5) < comb(23, 5) == 33649
    monomials = gradient_monomials(blocks)
    assert len(monomials) == 23 and all(len(row) == 77 for row in monomials)
    assert sum(0 in h for h in blocks) == 77
    assert sum(0 in h and 1 in h for h in blocks) == 21
    assert sum(0 in h and 1 in h and 2 in h for h in blocks) == 5
    return wt, B5, monomials


def main():
    parser = argparse.ArgumentParser()
    parser.add_argument('--json', type=Path, help='Write machine-readable heptads and invariant support')
    args = parser.parse_args()
    blocks = heptads()
    weights, B5, dJ7 = verify(blocks)
    print('Golay code G23: [23,12,7]')
    print('Weight enumerator:', dict(sorted(weights.items())))
    print('Steiner heptads:', len(blocks))
    print('Every 4-subset in exactly one heptad: verified')
    print('Distinguished 5-subsets (J5 terms):', len(B5))
    print('Heptad invariant J7 terms:', len(blocks))
    print('Gradient J7 monomials per coordinate:', len(dJ7[0]))
    if args.json:
        obj = {
            'coordinate_index_base': 0,
            'generator_polynomial_hex': hex(GOLAY_G),
            'heptads': blocks,
            'five_subsets_in_heptads': B5,
            'gradient_J7_monomials': dJ7,
            'interpretation': {
                'J5': 'sum of product(X_i) over five_subsets_in_heptads',
                'J7': 'sum of product(X_i) over heptads',
                'D_i': 'sum of product(X_j) over gradient_J7_monomials[i]',
                'beta_i': 'alpha_i + u*D_i(alpha)',
                'caveat': 'alpha must be labelled with the same M23 action; this does not compute a new rational splitting field by itself',
            },
        }
        args.json.write_text(json.dumps(obj, separators=(',', ':'))+'\n', encoding='utf-8')
        print('Wrote', args.json)


if __name__ == '__main__':
    main()
\end{lstlisting}

\section{Bibliography and source data}
\begin{thebibliography}{99}
\bibitem{HJLPPZ26} X. Huang, B. Jackson, K.-H. Lee, B. Poonen, R. Pries, and S. Zhang, \emph{The Mathieu group \(M_{23}\) is a Galois group over \(\QQ\)}, arXiv:2608.08538v2 (7 October 2026). \url{https://arxiv.org/abs/2608.08538}. Source files and computations: \url{https://github.com/shaowuz/m23isgalois}.
\bibitem{Janik26} J. A. Janik, \emph{Real Structure and the Galois Fixed Point in the Nielsen Class \((2,23A,23B)\) of \(M_{23}\)}, 7 October 2026 (supplied manuscript). Project archive: \url{https://codeberg.org/johnjanik/inverse_galois}.
\bibitem{MM18} G. Malle and B. H. Matzat, \emph{Inverse Galois Theory}, 2nd ed., Springer, 2018.
\bibitem{FV91} M. D. Fried and H. V\"olklein, \emph{The inverse Galois problem and rational points on moduli spaces}, Math. Ann. \textbf{290} (1991), 771--800.
\bibitem{Conway85} J. H. Conway et al., \emph{Atlas of Finite Groups}, Oxford University Press, 1985.
\end{thebibliography}
\end{document}
